\section{Taft Hopf algebras}\label{sec:blah}
We fix a prime $p$ and a primitive $p$th root of unity $\xi\in k$ in the following. The Taft Hopf algebra $T_p(\xi)$~\cite{taft} is generated as a $k$-algebra by $g$ and $x$ with the relations \begin{gather}
g^p=1,\qquad
x^p=0,
\\
gx=\xi xg.
\end{gather}
It is sometimes called the quantum ${\mathfrak{sl}}_2$ Borel algebra. It is $p^2$ dimensional over $k$. Furthermore, the coalgebra structure is
\begin{gather}
\Delta(g)=g\ot g,\qquad \Delta(x)=x\ot 1+g\ot x
\end{gather}
with $\epsilon(g)=1$, $\epsilon(x)=0$, and thus $S(g)=g^{-1}$, while $S(x)=-g^{-1}x$. Note that
\begin{gather}
S^2(x)=\xi^{-1}x\neq x,
\end{gather}
making $T_2(-1)$ the smallest Hopf algebra with $S^2\neq {\rm Id}$. The Taft algebra $T_2(-1)$ is somewhat different from the other $T_p(\xi)$ and has its own name: Sweedler’s Hopf algebra.

\subsection{The identification with the dual}
The Taft algebra is isomorphic to its dual. We~need some explicit formulas establishing the isomorphism $T_p(\xi)\simeq T_p(\xi)^*$ and its inverse. Suppose that $\omega$ is a $p$th root of unity, let
\begin{gather}
(n)_\omega=1+\cdots +\omega^{n-1}
\end{gather}
and
\begin{gather}
(n)_\omega !=(n)_\omega\cdots(1)_\omega.
\end{gather}

The verification of the following is left to the reader; key details can be found in~\cite{selfdual}. The lemma itself can be obtained from~\cite{nen}.
\begin{Lemma}\label{nenlem}
 As Hopf algebras
 \begin{gather}
 T_p(\xi)^*\simeq T_p(\xi).
 \end{gather}
\end{Lemma}
\begin{proof}
Consider a basis of $T_p(\xi)$: $\big\{g^ix^j\big\}_{ i,j=0}^{p-1}$ so that $\big\{\big(g^ix^j\big)^*\big\}$ denotes the dual basis of $T_p(\xi)^*$. Then the isomorphism of Hopf algebras and its inverse are given by
\begin{gather}
g^ix^j\mapsto(j)_{\xi^{-1}}!\sum_l\xi^{i(j+l)}\big(g^lx^j\big)^*
\end{gather}
and
\begin{gather}
(g^ix^j)^*\mapsto\dfrac{1}{p(j)_{\xi^{-1}}!}\sum_l\xi^{-l(i+j)}g^lx^j.
\end{gather}
\end{proof}

\begin{Corollary}\label{gens}The twisted double $\widehat{D}(T_p(\xi))$ is a quotient of $k\left\langle x,x',g,g'\right\rangle.$

\begin{itemize}\itemsep=0pt
\item The relations are
\begin{gather*}
x^p=x'^p=g^p-1=g'^p-1=0,
\\
gg'=g'g,\qquad
gx=\xi xg,\qquad
g'x'=\xi x'g',\qquad
gx'=\xi^{-1} x'g, \qquad
g'x=\xi^{-1} xg',
\\
xx'-\xi^{-1}x'x=1-\xi^{-1}g'^{-1}g.
\end{gather*}

\item The actions of $g'$ and $g$ on a $\widehat{D}(T_p(\xi))$-module $V$, yields a $(\mathbb{Z}/p)^2$-grading on $V$ by their eigenspaces, i.e., $g'$, $g$ act on $V_{ij}$ by $\xi^i$, $\xi^j$, respectively. Thus $x$ and $x'$ have degrees $(-1,1)$ and $(1,-1)$, respectively.

\item The $S^1$-action of $\sigma$ on $V_{ij}$ is
\begin{gather}\label{sigmaaction}
\sum_{l=0}^{p-1}\dfrac{\xi^{(i-l)(j+l)}}{(l)_{\xi^{-1}}!}x'^lx^l.
\end{gather}
\end{itemize}
\end{Corollary}

\begin{proof}
We use the identification of vector spaces $\widehat{D}(T_p(\xi))=(T_p(\xi))^*\ot T_p(\xi)$ as in Lemma~\ref{bas} followed by $(T_p(\xi))^*\ot T_p(\xi) \simeq T_p(\xi)\ot T_p(\xi)$ from Lemma~\ref{nenlem}. We~let
\begin{gather}
x'=x\ot 1, \qquad
x=1\ot x, \qquad
g'=g\ot 1, \qquad
g=1\ot g
\end{gather}
in the latter. To derive the rest of the relations we apply \eqref{mult1}. The action of $\sigma=\sum_{ij}\big(g^ix^j\big)^*\ot g^ix^j$ is computed on the graded components directly.
\end{proof}

Observe that it follows from Corollary~\ref{gens} that $gg'\in\widehat{D}(T_p(\xi))$ is central. Since we see that $(gg')^p=1$ so its action on $\widehat{D}(T_p(\xi))$ is diagonalizable with eigenvalues $\xi^s$, $s\in\zp$. Thus as an~algebra
\begin{gather}\label{split}
\widehat{D}(T_p(\xi))=\bigoplus_s \widehat{D}(T_p(\xi))/(gg'-\xi^s)
\end{gather}
so that it suffices to understand $\widehat{D}(T_p(\xi))/(gg'-\xi^s)$. There are two cases: $p=2$ and $p>2$. We~will begin by briefly discussing the latter (though without addressing the $S^1$-action), and then concentrate our attention on the former (with examining the $S^1$-action) to achieve the goal set out in the abstract.

